\documentclass[11pt]{article}
\usepackage[T1]{fontenc}
\usepackage{lmodern,geometry,amsmath,amssymb,microtype,xcolor}
\usepackage[colorlinks=true,urlcolor=blue]{hyperref}
\geometry{letterpaper,left=.87in,right=.87in,top=.85in,bottom=.77in}
\definecolor{ink}{HTML}{172D45}
\definecolor{teal}{HTML}{146565}
\setlength{\parindent}{0pt}
\setlength{\parskip}{5pt}
\pagestyle{empty}
\newcommand{\entry}[2]{\vspace{3pt}{\large\sffamily\bfseries\color{ink}#1}\par #2\par}
\begin{document}
{\sffamily\color{teal} MATHEMATICAL RESEARCH NOTES\hfill 25 NOVEMBER 2026}\par
\vspace{14pt}
{\fontsize{27}{31}\selectfont\sffamily\bfseries\color{ink}
Beyond Dirichlet GRH}\par
{\large Heavenly elliptic curves, critical-line Euler products, and geometric progressions}\par
\vspace{3pt}{\color{teal}\hrule}\vspace{5pt}
\entry{1. Both cubic finiteness questions have affirmative answers}{
Suzanne O'Hara asks whether only finitely many heavenly elliptic curves and field--prime pairs arise over all Galois cubic fields at primes $\ell\ge11$, after excluding geometric classes already heavenly over $\mathbb Q$ at the same prime. Rasmussen--Tamagawa bound $\ell$ uniformly in the cubic field. McLeman--Rasmussen bound the number of geometric classes for each such $\ell$, and O'Hara bounds each fiber over a geometric class. This qualitative answer is unconditional. The same exclusion suffices in odd prime degree $d$ for $\ell>2d+1$, and for the $A_5$ sextic fields at $\ell\ge17$, using Dirichlet GRH.}
\entry{2. One bound for each of two infinite families of sextic fields}{
For every dimension $g$, the uniform Rasmussen--Tamagawa conjecture holds over sextic fields with normal-closure group $A_4$ or $A_5$. The argument takes the odd residue degrees in the earlier ProveIt report and applies them to the Frobenius obstruction for heavenly abelian varieties. At $g=1$, only 29 primes survive, all at most $4{,}027$. The same list works for all cubic fields, leaving 25 primes in O'Hara's range. The proof includes ramified primes, a finite sieve and an analytic bound for the infinite tail. Attainment of the maximum is undetermined.}
\entry{3. A critical-line Euler product that converges at every cutoff}{
Let $P_x(s,\chi)=\prod_{p\le x}(1-\chi(p)p^{-s})^{-1}$. Two offsets with opposite phases at the cutoff cancel the common sharp prime-sum error. For the character $\chi_4$ this gives $|P_x(1/2+\pi i/(2\log x),\chi_4)|\to0.547115431519\ldots$. A separate inverse theorem recovers the progression modulus and quadratic twist from the signed, aggregated zero multiplicities in the Lean proof. The ordinary fixed-center deep Riemann hypothesis remains open.}
\entry{4. Which geometric progressions can occur in a periodic set?}{
The repository's weighted Behrend construction yields dense sets of residues with prescribed geometric-progression ratios and quantitative control of their gaps. The result gives a finite-modulus restriction for the bounded-gap conjecture.}
\vfill
{\footnotesize\sffamily\color{ink} The archive includes the four papers, LaTeX sources, exact arithmetic checkers and a numerical Euler-product illustration. The downstream theorems are paper proofs. Upstream scope and independent verification are documented in the analytic paper and the README.}
\end{document}
